\documentclass[11pt]{article}

% CPSC 250 lecture-note style handout
\usepackage[margin=0.85in]{geometry}
\usepackage[T1]{fontenc}
\usepackage[utf8]{inputenc}
\usepackage{lmodern}
\usepackage{microtype}
\usepackage{xcolor}
\usepackage{hyperref}
\usepackage{enumitem}
\usepackage{booktabs}
\usepackage{tabularx}
\usepackage{fancyhdr}
\usepackage{titlesec}
\usepackage{tcolorbox}
\usepackage{listings}
\usepackage{multicol}
\usepackage{amsmath}

\definecolor{cnuBlue}{HTML}{003A70}
\definecolor{cnuSilver}{HTML}{A2AAAD}
\definecolor{codeBack}{HTML}{F6F8FA}
\definecolor{codeFrame}{HTML}{D0D7DE}
\definecolor{noteBack}{HTML}{EEF6FF}
\definecolor{warnBack}{HTML}{FFF8E5}

\hypersetup{
  colorlinks=true,
  linkcolor=cnuBlue,
  urlcolor=cnuBlue
}

\setlength{\headheight}{14pt}
\pagestyle{fancy}
\fancyhf{}
\lhead{CPSC 250 / 250L}
\chead{Git Command Line Summary}
\rhead{Summer Course Reference}
\cfoot{\thepage}
\renewcommand{\headrulewidth}{0.4pt}

\titleformat{\section}{\Large\bfseries\color{cnuBlue}}{\thesection}{0.75em}{}
\titleformat{\subsection}{\large\bfseries\color{cnuBlue}}{\thesubsection}{0.75em}{}
\titleformat{\subsubsection}{\bfseries\color{cnuBlue}}{\thesubsubsection}{0.75em}{}

\lstdefinestyle{terminal}{
  basicstyle=\ttfamily\small,
  backgroundcolor=\color{codeBack},
  frame=single,
  rulecolor=\color{codeFrame},
  columns=fullflexible,
  keepspaces=true,
  showstringspaces=false,
  breaklines=true,
  upquote=true
}

\newtcbox{\gitcmd}{on line, boxrule=0.4pt, arc=2pt, colback=codeBack, colframe=codeFrame,
  left=2pt, right=2pt, top=1pt, bottom=1pt, boxsep=1pt, fontupper=\ttfamily\small}

\newtcolorbox{lecturebox}[1]{
  colback=noteBack,
  colframe=cnuBlue,
  title=\textbf{#1},
  arc=2pt,
  boxrule=0.6pt
}

\newtcolorbox{warningbox}[1]{
  colback=warnBack,
  colframe=orange!70!black,
  title=\textbf{#1},
  arc=2pt,
  boxrule=0.6pt
}

\newcommand{\term}[1]{\textbf{\color{cnuBlue}#1}}
\newcommand{\repo}{\texttt{cpsc250L}}
\newcommand{\main}{\texttt{main}}
\newcommand{\origin}{\texttt{origin}}
\newcommand{\upstream}{\texttt{upstream}}

\title{\vspace{-0.5in}\textbf{CPSC 250 / 250L\\Command Line Git Reference}\\
\large A lecture-style summary with examples}
\author{}
\date{}

\begin{document}
\maketitle
\thispagestyle{fancy}

\begin{lecturebox}{Big idea}
Git is a version-control system.  It keeps a history of your work, lets you move safely between versions, and connects your local copy of a project to a shared copy on GitHub.  In this course, the usual workflow is:
\[
\text{edit files} \longrightarrow \text{stage files} \longrightarrow \text{commit locally} \longrightarrow \text{push to GitHub}.
\]
\end{lecturebox}

\tableofcontents
\newpage

\section{The mental model}

A Git project has several important places where files can live.

\begin{center}
\begin{tabularx}{0.95\textwidth}{@{}lX@{}}
\toprule
\term{Place} & \term{Meaning} \\
\midrule
Working directory & The actual files you are editing in PyCharm, VS Code, or a terminal. \\
Staging area & A preparation area for the next commit.  Files get here with \gitcmd{git add}. \\
Local repository & The saved history on your own computer.  New snapshots get here with \gitcmd{git commit}. \\
Remote repository & The GitHub copy of the repository.  Your fork is usually called \origin.  The instructor repository is usually called \upstream. \\
\bottomrule
\end{tabularx}
\end{center}

\begin{lecturebox}{Course vocabulary}
\begin{itemize}[itemsep=2pt]
  \item \term{Repository}: a folder tracked by Git.
  \item \term{Commit}: a saved snapshot of the project.
  \item \term{Branch}: a named line of development.
  \item \term{Fork}: your GitHub copy of someone else's repository.
  \item \term{Remote}: a nickname for a GitHub repository connected to your local repository.
  \item \term{Merge}: combining changes from one branch or repository into another.
\end{itemize}
\end{lecturebox}

\section{One-time setup commands}

These commands are usually done once on a computer or once for a repository.

\subsection{Configure your name and email}
Git records the author of each commit.  The name and email should identify you.

\begin{lstlisting}[style=terminal]
git config --global user.name "Ada Lovelace"
git config --global user.email "ada.lovelace.24@cnu.edu"
\end{lstlisting}

Check your settings:

\begin{lstlisting}[style=terminal]
git config --global --list
\end{lstlisting}

\subsection{Clone a repository}
\gitcmd{git clone} downloads a repository from GitHub and creates a local working copy.

\begin{lstlisting}[style=terminal]
git clone https://github.com/YOUR-USERNAME/cpsc250L.git
\end{lstlisting}

Then move into the folder:

\begin{lstlisting}[style=terminal]
cd cpsc250L
\end{lstlisting}

\subsection{See the connected remotes}
\gitcmd{git remote -v} shows the GitHub repositories connected to your local repository.

\begin{lstlisting}[style=terminal]
git remote -v
\end{lstlisting}

Typical course output:

\begin{lstlisting}[style=terminal]
origin    https://github.com/YOUR-USERNAME/cpsc250L.git (fetch)
origin    https://github.com/YOUR-USERNAME/cpsc250L.git (push)
upstream  https://github.com/brash99/cpsc250L.git (fetch)
upstream  https://github.com/brash99/cpsc250L.git (push)
\end{lstlisting}

\subsection{Add the instructor repository as upstream}
If your local repository only knows about your fork, add the instructor repository as \upstream.

\begin{lstlisting}[style=terminal]
git remote add upstream https://github.com/brash99/cpsc250L.git
\end{lstlisting}

\subsection{Change a remote URL}
Sometimes a remote points to the wrong GitHub repository.  Change it with \gitcmd{git remote set-url}.

\begin{lstlisting}[style=terminal]
git remote set-url upstream https://github.com/brash99/cpsc250L-instructor.git
\end{lstlisting}

Then verify:

\begin{lstlisting}[style=terminal]
git remote -v
\end{lstlisting}

\section{The daily save-and-submit workflow}

\subsection{Check the current state}
The most important Git command is \gitcmd{git status}.  Use it constantly.

\begin{lstlisting}[style=terminal]
git status
\end{lstlisting}

Common messages:

\begin{center}
\begin{tabularx}{0.95\textwidth}{@{}lX@{}}
\toprule
\term{Message} & \term{Meaning} \\
\midrule
\texttt{nothing to commit, working tree clean} & Git sees no unsaved changes. \\
\texttt{Changes not staged for commit} & Files changed, but they are not yet prepared for the next commit. \\
\texttt{Untracked files} & Git sees new files that have not yet been added. \\
\texttt{Changes to be committed} & Files are staged and ready to commit. \\
\bottomrule
\end{tabularx}
\end{center}

\subsection{Stage files}
To stage one file:

\begin{lstlisting}[style=terminal]
git add labs/lab03_functions/solution.py
\end{lstlisting}

To stage every changed file in the current repository:

\begin{lstlisting}[style=terminal]
git add .
\end{lstlisting}

\begin{warningbox}{Important habit}
Before using \gitcmd{git add .}, run \gitcmd{git status}.  Make sure you are not accidentally adding temporary files, large data files, or virtual-environment folders.
\end{warningbox}

\subsection{Commit staged files}
A commit saves a snapshot in your local repository.

\begin{lstlisting}[style=terminal]
git commit -m "Complete lab 03 functions"
\end{lstlisting}

Good commit messages are short, specific, and written in the imperative style:

\begin{lstlisting}[style=terminal]
git commit -m "Add CSV parsing for lab 04"
git commit -m "Fix Plant __str__ method"
git commit -m "Refactor book search into helper function"
\end{lstlisting}

\subsection{Push your commits to GitHub}
A commit is still only on your computer until you push it.

\begin{lstlisting}[style=terminal]
git push
\end{lstlisting}

The first push from a new branch may need:

\begin{lstlisting}[style=terminal]
git push -u origin lab05-feature
\end{lstlisting}

After that, \gitcmd{git push} is usually enough.

\section{Getting updates from GitHub}

\subsection{Fetch without changing your files}
\gitcmd{git fetch} contacts GitHub and downloads information about new commits, but it does not merge them into your current branch.

\begin{lstlisting}[style=terminal]
git fetch upstream
\end{lstlisting}

This is a safe way to ask, ``What has changed in the instructor repository?''

\subsection{Pull from your own GitHub fork}
\gitcmd{git pull} fetches and then merges changes into your current branch.

\begin{lstlisting}[style=terminal]
git pull
\end{lstlisting}

This is commonly used when your GitHub fork has changes that your local computer does not yet have.

\subsection{Update your local main branch from the instructor repository}
A common course pattern is:

\begin{lstlisting}[style=terminal]
git switch main
git fetch upstream
git merge upstream/main
\end{lstlisting}

This says:

\begin{enumerate}[itemsep=2pt]
  \item Move to your local \main{} branch.
  \item Get the latest information from the instructor repository.
  \item Merge the instructor's \main{} branch into your local \main{} branch.
\end{enumerate}

Then push the updated version of your fork:

\begin{lstlisting}[style=terminal]
git push origin main
\end{lstlisting}

\section{Looking at history and changes}

\subsection{View the commit history}

\begin{lstlisting}[style=terminal]
git log
\end{lstlisting}

A compact version is often easier to read:

\begin{lstlisting}[style=terminal]
git log --oneline
\end{lstlisting}

A useful branch graph view:

\begin{lstlisting}[style=terminal]
git log --oneline --graph --decorate --all
\end{lstlisting}

\subsection{See unstaged changes}
\gitcmd{git diff} shows changes in the working directory that have not been staged.

\begin{lstlisting}[style=terminal]
git diff
\end{lstlisting}

\subsection{See staged changes}
After \gitcmd{git add}, use:

\begin{lstlisting}[style=terminal]
git diff --staged
\end{lstlisting}

This answers: ``What exactly am I about to commit?''

\section{Branches}

Branches allow us to work on a feature or lab without immediately changing \main{}.

\subsection{List branches}

\begin{lstlisting}[style=terminal]
git branch
\end{lstlisting}

Show local and remote branches:

\begin{lstlisting}[style=terminal]
git branch -a
\end{lstlisting}

\subsection{Create and switch to a new branch}

\begin{lstlisting}[style=terminal]
git switch -c lab05-feature
\end{lstlisting}

Older Git examples often use \gitcmd{checkout} instead:

\begin{lstlisting}[style=terminal]
git checkout -b lab05-feature
\end{lstlisting}

\subsection{Switch between existing branches}

\begin{lstlisting}[style=terminal]
git switch main
git switch lab05-feature
\end{lstlisting}

Older style:

\begin{lstlisting}[style=terminal]
git checkout main
\end{lstlisting}

\subsection{Merge a feature branch into main}

\begin{lstlisting}[style=terminal]
git switch main
git merge lab05-feature
\end{lstlisting}

Then push:

\begin{lstlisting}[style=terminal]
git push origin main
\end{lstlisting}

\subsection{Delete a finished branch}
After the branch has been merged:

\begin{lstlisting}[style=terminal]
git branch -d lab05-feature
\end{lstlisting}

Force-delete a branch only when you are sure you do not need it:

\begin{lstlisting}[style=terminal]
git branch -D lab05-feature
\end{lstlisting}

\section{Undoing and fixing mistakes}

\subsection{Unstage a file}
If you staged a file with \gitcmd{git add} but do not want it in the next commit:

\begin{lstlisting}[style=terminal]
git restore --staged labs/lab04_csv/solution.py
\end{lstlisting}

Older Git examples may use:

\begin{lstlisting}[style=terminal]
git reset HEAD labs/lab04_csv/solution.py
\end{lstlisting}

\subsection{Discard local changes to a file}
This throws away uncommitted changes to a file and restores the last committed version.

\begin{lstlisting}[style=terminal]
git restore labs/lab04_csv/solution.py
\end{lstlisting}

\begin{warningbox}{This is destructive}
Only use \gitcmd{git restore FILE} when you really want to lose the uncommitted edits in that file.
\end{warningbox}

\subsection{Fix the most recent commit message}

\begin{lstlisting}[style=terminal]
git commit --amend -m "Complete lab 04 CSV processing"
\end{lstlisting}

Use this before pushing when possible.  Amending a commit that has already been pushed can confuse shared history.

\subsection{Revert a committed change safely}
\gitcmd{git revert} creates a new commit that undoes an earlier commit.

\begin{lstlisting}[style=terminal]
git log --oneline
git revert abc1234
\end{lstlisting}

This is safer than rewriting history because it preserves the record of what happened.

\section{Merge conflicts}

A merge conflict happens when Git cannot automatically combine two versions of a file.

\subsection{Typical conflict workflow}

\begin{lstlisting}[style=terminal]
git status
# Open the conflicted file and edit it by hand.
git add conflicted_file.py
git commit
\end{lstlisting}

Conflict markers look like this:

\begin{lstlisting}[style=terminal]
<<<<<<< HEAD
print("My version")
=======
print("Other version")
>>>>>>> upstream/main
\end{lstlisting}

Your job is to edit the file into the version that should remain, then remove the conflict markers.

\section{Files Git should ignore}

Some files do not belong in the repository: virtual environments, cache files, temporary files, and generated output.  Put patterns for these files in \texttt{.gitignore}.

Example \texttt{.gitignore} entries for Python work:

\begin{lstlisting}[style=terminal]
.venv/
venv/
__pycache__/
*.pyc
.DS_Store
.idea/
\end{lstlisting}

After editing \texttt{.gitignore}, commit it like any other file:

\begin{lstlisting}[style=terminal]
git add .gitignore
git commit -m "Update Python gitignore rules"
git push
\end{lstlisting}

\section{Common course scenarios}

\subsection{Scenario 1: Submit today's lab}

\begin{lstlisting}[style=terminal]
git status
git add labs/lab04_csv_processing/solution.py
git status
git commit -m "Complete lab 04 CSV processing"
git push
\end{lstlisting}

\subsection{Scenario 2: Get the latest lab files from the instructor repository}

\begin{lstlisting}[style=terminal]
git switch main
git status
git fetch upstream
git merge upstream/main
git push origin main
\end{lstlisting}

\subsection{Scenario 3: Start a new lab branch}

\begin{lstlisting}[style=terminal]
git switch main
git pull
git switch -c lab05-classes
\end{lstlisting}

Work on the lab, then save it:

\begin{lstlisting}[style=terminal]
git add labs/lab05_classes_feature_branches/
git commit -m "Complete lab 05 classes"
git push -u origin lab05-classes
\end{lstlisting}

\subsection{Scenario 4: Merge your lab branch back into main}

\begin{lstlisting}[style=terminal]
git switch main
git merge lab05-classes
git push origin main
\end{lstlisting}

\subsection{Scenario 5: Diagnose a weird repository problem}
Start with read-only commands:

\begin{lstlisting}[style=terminal]
git status
git branch -a
git remote -v
git log --oneline --graph --decorate --all
\end{lstlisting}

These commands do not change the repository.  They help you see where you are before deciding what to do.

\section{Command summary table}

\begin{center}
\small
\begin{tabularx}{\textwidth}{@{}lX@{}}
\toprule
\term{Command} & \term{Purpose} \\
\midrule
\gitcmd{git config --global user.name "Name"} & Set the author name for commits. \\
\gitcmd{git config --global user.email "email"} & Set the author email for commits. \\
\gitcmd{git clone URL} & Download a repository from GitHub. \\
\gitcmd{git status} & Show changed, staged, and untracked files. \\
\gitcmd{git add FILE} & Stage a file for the next commit. \\
\gitcmd{git add .} & Stage all changes in the current repository. \\
\gitcmd{git commit -m "message"} & Save a local snapshot. \\
\gitcmd{git push} & Upload local commits to GitHub. \\
\gitcmd{git pull} & Download and merge changes from the default remote branch. \\
\gitcmd{git fetch upstream} & Download information from the instructor repository without merging. \\
\gitcmd{git merge upstream/main} & Merge instructor changes into the current branch. \\
\gitcmd{git remote -v} & List connected GitHub repositories. \\
\gitcmd{git remote add upstream URL} & Add the instructor repository as a remote. \\
\gitcmd{git remote set-url upstream URL} & Change where a remote points. \\
\gitcmd{git branch} & List local branches. \\
\gitcmd{git branch -a} & List local and remote branches. \\
\gitcmd{git switch BRANCH} & Move to an existing branch. \\
\gitcmd{git switch -c BRANCH} & Create and move to a new branch. \\
\gitcmd{git checkout BRANCH} & Older command for switching branches. \\
\gitcmd{git checkout -b BRANCH} & Older command for creating and switching to a branch. \\
\gitcmd{git log --oneline} & Show compact commit history. \\
\gitcmd{git diff} & Show unstaged changes. \\
\gitcmd{git diff --staged} & Show staged changes. \\
\gitcmd{git restore --staged FILE} & Unstage a file. \\
\gitcmd{git restore FILE} & Discard uncommitted changes to a file. \\
\gitcmd{git commit --amend} & Modify the most recent commit. \\
\gitcmd{git revert HASH} & Safely undo an old commit by creating a new commit. \\
\gitcmd{git branch -d BRANCH} & Delete a merged local branch. \\
\bottomrule
\end{tabularx}
\end{center}

\section{A final checklist before asking for help}

When Git seems confusing, run these commands and read the output carefully:

\begin{lstlisting}[style=terminal]
git status
git branch -a
git remote -v
git log --oneline --graph --decorate --all
\end{lstlisting}

Then ask:

\begin{enumerate}[itemsep=2pt]
  \item What branch am I on?
  \item Are my files changed, staged, committed, or pushed?
  \item Is \origin{} my fork?
  \item Is \upstream{} the instructor repository?
  \item Am I trying to save my work, submit my work, or get new instructor files?
\end{enumerate}

\begin{lecturebox}{The safest Git habit}
Use \gitcmd{git status} before and after every important Git command.  Git usually tells you exactly what state the repository is in and what command might make sense next.
\end{lecturebox}

\end{document}
